\documentclass[10pt,a4paper]{amsart}
\usepackage[margin=19mm]{geometry}
\usepackage{amsmath,amssymb,mathtools}
\usepackage[scr=rsfs,cal=euler]{mathalfa}
\usepackage{xfrac}
\usepackage[unicode,hidelinks,bookmarks=false]{hyperref}
\newcommand{\ie}{i.e.\ }
\newcommand{\cf}{cf.\ }
\newcommand{\resp}{respectively\ }
\newcommand{\Subsection}{\subsection}
\providecommand{\nobreakdash}{\nobreak\hbox{-}}
\setlength{\emergencystretch}{2em}
\multlinegap0pt
\usepackage{amssymb}
\newtheorem{lem}{Lemma}
\title{A proof of Hare's $\epsilon$ unfair conjecture}
\author{Tho Nguyen Xuan, Duc Hiep Pham, Huong Pham Lan}
\date{Source: arXiv:2609.09215v1 (CC BY 4.0)}
\begin{document}
\maketitle

\noindent\emph{Source and licence.} A proof of Hare's $\epsilon$ unfair conjecture. Version arXiv:2609.09215v1 (CC BY 4.0), distributed by arXiv under the Creative Commons Attribution 4.0 International licence, http://creativecommons.org/licenses/by/4.0/, as recorded in the arXiv licence metadata for this exact version and shown on the abstract page https://arxiv.org/abs/2609.09215v1. Full licence notice: \texttt{licenses/arxiv-2609.09215-CC-BY-4.0.txt}.\par\smallskip

\noindent\textit{Editorial scope.} Counted unit: the article's lem L1 (source0.tex lines 113--119). The article's complete original proof of it is delivered with it (source0.tex lines 120--159, one \texttt{proof} environment of 40 lines). No mathematics is written, repaired or re-derived: the statement and the proof are the article's own lines, in the article's own notation, and no retained line is substituted or re-flowed.  Retained uncounted as the source-local context the proof needs: lines 100--102; lines 104--110.  Nothing was omitted from any retained span, so the package carries no omission record.  No retained line contains a citation, so the wrapper carries no bibliography and no bibliography entry.  No retained line contains a figure environment, an image-inclusion command or drawing code, so no image is delivered and the package declares no asset dependency.  Editorial formatting only, in addition: an AMS \texttt{amsart} preamble that restates the article's own declarations and macros used by the retained lines, namely the environments \texttt{lem} on separate counters, together with the standard packages those lines need.  The retained environments print in this document as Lemma 1 in order of appearance, whereas the article prints the same items with the numbering of its own full document.  The complete native source of the article as distributed in the arXiv e-print for this exact version is bundled unchanged as \texttt{sources/209791/source0.tex}.
Since $m\pi/3=(2n+1)\pi$, we have for all real $\delta$ then \begin{equation}\label{E6}
G(\dfrac{\pi}{3}+\delta)=1-2\cos(m\delta)+\cos\delta-\sqrt{3}\sin\delta.
\end{equation}

\begin{equation}\label{ineq}
\begin{cases}
1-\dfrac{y^2}{2}\leq\cos y\leq 1-\dfrac{y^2}{2}+\dfrac{y^4}{24},\\
y-\dfrac{y^3}{6}\leq \sin y\leq y,\\
\cos y\geq 1-\dfrac{y^2}{2}+\dfrac{y^4}{24}-\dfrac{y^6}{720}.
\end{cases}
\end{equation}

\begin{lem}\label{L1}
For every $n\geq 1$, there exists $\delta^{*}\in \left(\dfrac{\sqrt{3}}{m^2},\dfrac{\sqrt{3}}{m^2}+\dfrac{3}{m^4}\right)$
such that $G(\pi/3+\delta^*)=0$. Let $\theta^*=\pi/3+\delta^*$ and $t:=1-2\cos\theta^*$. Then \begin{equation}\label{E8}
\dfrac{3}{m^2}-\dfrac{2}{m^6}\leq t\leq \dfrac{3}{m^2}+\dfrac{7}{m^4}<\dfrac{3.1}{m^2}<1.
\end{equation}
Moreover, $\delta^*<1.78m^{-2}$ and $\theta^*\in (\pi/3,\pi/2)$.
\end{lem}

\begin{proof}
Let $\delta_1=\sqrt{3}/m^2$ and $\delta_2=\sqrt{3}/m^2+3/m^4$. Let $y_1=m\delta_1$ and $y_2=m\delta_2$. For $m\geq 9$, \[\delta_2=\dfrac{\sqrt{3}}{m^2}+\dfrac{3}{m^4}\leq \dfrac{\sqrt{3}}{m^2}+\dfrac{1}{27m^2}<\dfrac{1.78}{m^2},\]
\[y_2=m\delta_2<\dfrac{1.78}{m}.\]
From  \eqref{E6} and \eqref{ineq}, we have 
\[\begin{split}
G\left(\dfrac{\pi}{3}+\delta_1\right)&\leq \left[-1+y_1^2-\dfrac{y_1^4}{12}+\dfrac{y_1^6}{360}\right] +\left[1-\dfrac{\delta_1^2}{2}+\dfrac{\delta_1^4}{24}\right]+\left[-\sqrt{3}\delta_1+\dfrac{\sqrt{3}\delta_1^3}{6}\right]\\
&=\left(
-1+\frac{3}{m^2}-\frac{3}{4m^4}+\frac{3}{40m^6}
\right)+\left(
1-\frac{3}{2m^4}+\frac{3}{8m^8}
\right)+\left(
-\frac{3}{m^2}+\frac{3}{2m^6}
\right)\\
&=\dfrac{-90m^4+63m^2+15}{40m^8}\\
&<0,
\end{split}\]
and
\[\begin{split}
G\left(\dfrac{\pi}{3}+\delta_2\right)&\geq \left[-1+y_2^2-\dfrac{y_2^4}{12}\right]+\left[1-\dfrac{\delta_2^2}{2}\right]-\sqrt{3}\delta_2\\
&=y_2^2-\sqrt{3}\delta_2-\dfrac{y_2^4}{12}-\dfrac{\delta_2^2}{2},
\end{split}\]
Note that \[y_2^2-\sqrt{3}\delta_2=\dfrac{3\sqrt{3}}{m^4}+\dfrac{9}{m^6}>\dfrac{3\sqrt{3}}{m^4}.\]
Hence,
\[G\left(\dfrac{\pi}{3}+\delta_2\right)>\dfrac{3\sqrt{3}-\dfrac{1.78^4}{12}-\dfrac{1.78^2}{2}}{m^4}>0.\]
By continuity, there exists $\delta^*\in (\delta_1,\delta_2)$ such that $G(\pi/3+\delta^*)=0$.
Finally, \[t=1-2\cos(\dfrac{\pi}{3}+\delta^*)=1-\cos\delta^*+\sqrt{3}\sin\delta^*>0.\]
By \eqref{ineq}, we also have
\[\begin{split}t&\leq \dfrac{(\delta^*)^2}{2}+\sqrt{3}\delta^*\\
&<\dfrac{\delta_2^2}{2}+\sqrt{3}\delta_2\\
&<\dfrac{1}{2}\left(\dfrac{1.78}{m^2}\right)^2+\dfrac{3}{m^2}+\dfrac{3\sqrt{3}}{m^4}\\
&<\dfrac{3}{m^2}+\dfrac{7}{m^4},\end{split}\]
and \[\begin{split}
t&\geq \sqrt{3}\sin\delta^*\\
&\geq \sqrt{3}\left(\delta^*-\dfrac{(\delta^*)^3}{6}\right)\\
&>\sqrt{3}\left(\delta_1-\dfrac{\delta_2^3}{6}\right)\\
&>\dfrac{3}{m^2}-\dfrac{\sqrt{3}\times 1.78^3}{6m^6}\\
&>\dfrac{3}{m^2}-\dfrac{2}{m^6}.
\end{split}\]
Therefore, \eqref{E8} is proved. For the last assertions, $\delta^*<\delta_2<1.78m^{-2}$ and $t\in(0,1)$ gives $\theta^*\in (\pi/3,\pi/2)$.
\end{proof}

\end{document}
